\documentclass[11pt]{article}
\usepackage{fontspec,amsmath,amssymb,amsthm,hyperref}
\usepackage[margin=1in]{geometry}
\setlength{\emergencystretch}{3em}
\title{Sharp initialization radius from actual transport}
\author{LangevinLean mathematical exposition}
\date{}
\begin{document}\maketitle
\noindent Source module: \texttt{AverageLMC/InitialMoment.lean}. Compile this proof document with XeLaTeX or LuaLaTeX. The accompanying HTML pairs the mathematical argument with the exact Lean statements and records the scope of the formalization.


The centered stochastic-gradient proof needs the initial squared distance from the true minimizer. This theorem derives it from finite initial transport and the sharp Gibbs concentration radius, rather than assuming a separate initialization moment bound.



For any true minimizer \(x_\star\), strong convexity gives the sharp Gibbs radius, and the coupling argument bounds the initial radius by the actual transport distance:


\[\int\|x-x_\star\|^2\,d\pi(x)\le\frac pm.\]
\[\int\|x-x_\star\|^2\,d\nu(x)\le2W_2^2(\nu,\pi)+\frac{2p}{m}.\]
\section{Integrate the two-term radius inequality and take the infimum}


Choose any finite-cost coupling of \ensuremath{\mu} and the actual Gibbs law. Both marginals have finite second moments: the Gibbs moment follows from strong convexity, and the \ensuremath{\mu} moment follows from finite transport. At each pair (x,y), write x−anchor=(x−y)+(y−anchor); the squared norm is at most twice the squared coupling displacement plus twice the squared target radius. Integrate, use the two exact marginal pushforwards, and identify the coupling displacement integral with the square of its L\ensuremath{{}^{2}} coupling cost. The true minimizer and strong convexity give the sharp target radius at most p/m. Thus every finite coupling bounds the initial radius by twice its squared cost plus 2p/m. Passing to the actual transport infimum proves the asserted inequality without assuming an optimal coupling or a Hessian regularity condition.



\noindent Lean declarations: \texttt{initial\_radius\_le}.

\end{document}
